% !PRINTED 703
\weberpage{703}
%<<CONT
der Lehre von den unendlichen Reihen eine Summe von Gliedern
der Form
\[
\left(\frac{1}{N(\mathfrak{p})^k}\,\frac{1}{N(\mathfrak{p}')^{k'}}\,\frac{1}{N(\mathfrak{p}'')^{k''}}\cdots\right)^s = \left(\frac{1}{N(\mathfrak{p}^k\mathfrak{p}'^{k'}\mathfrak{p}''^{k''}\dots)}\right)^s,
\]
die alle in der Form $N(\mathfrak{a})^{-s}$ enthalten sind, und jedes solche
Glied ergiebt sich ein- und nur einmal. Danach ist also
\weq{5}{\varPhi(s) = \prod \frac{1}{1-N(\mathfrak{p})^{-s}}.}

Sind daher $\mathfrak{p}_1, \mathfrak{p}_2, \dots, \mathfrak{p}_e$ die von einander verschiedenen
Primfactoren von $p$, und $f_1, f_2, \dots, f_e$ ihre Grade (\S.\,143), so
ergiebt sich
\weq{6}{\varPhi(s) = \overset{p}{\prod} \frac{1}{(1-p^{-sf_1})(1-p^{-sf_2})\dots(1-p^{-sf_e})},}
worin das unendliche Product $\prod$ über alle natürlichen Primzahlen
 $p$ auszudehnen ist, und nach (3), (4)
\weq{7}{gh = \mathop{\mathrm{Lim}}\limits_{s=1}\,(s-1)\varPhi(s).}

Hieraus lässt sich ein für mannigfache Anwendungen wichtiger
Schluss ziehen:

Wenn wir aus der Entwickelung
\[
1 + \frac{1}{p^s} + \frac{1}{p^{2s}} + \cdots = \frac{1}{1-p^{-s}}
\]
das Product für alle Primzahlen $p$ bilden, so ergiebt sich:
\weq{8}{\overset{p}{\prod}\frac{1}{1-p^{-s}} = \overset{n}{\sum}\frac{1}{n^s},}
worin sich die Summe auf der rechten Seite auf alle natürlichen
Zahlen $n$ erstreckt, und es hat also dies Product für alle Werthe
von $s$, die grösser als 1 sind, einen endlichen Werth. Der reciproke
Werth, nämlich das Product
\weq{9}{P = \overset{p}{\prod}(1-p^{-s}),}
dessen Factoren sämmtlich kleiner als 1 sind, hat daher, so
lange $s>1$ ist, einen von Null verschiedenen Werth. Diese
Eigenschaft bleibt nun erhalten, wenn wir bei der Bildung des
Productes $P$ nur einen Theil aller Primzahlen $p$ berücksichtigen,
weil das Product durch Weglassen beliebiger Factoren nur
vergrössert wird, ohne die Einheit je zu übersteigen.
